% docs/manual/developer/chapters/10-testing.tex
% 第 10 章：测试

\chapter{测试}

\openbench{} 的测试套件覆盖 config / data / runner / cli / manual generators 五大块。本章帮你看懂测试目录、跑测试、写新测试。

\section{测试目录布局}

\begin{minted}{text}
tests/
├── conftest.py             # 共享 fixtures（如有）
├── test_smoke.py           # 烟测：导入 + 版本
├── test_cli_integration.py # CLI 全部命令的集成测试
├── test_cli_stubs.py       # CLI 占位测试
├── test_dead_code_cleanup.py
├── test_processing_registry_cache.py
├── test_config/
│   ├── fixtures/                    # YAML / JSON fixture
│   │   ├── minimal.yaml
│   │   ├── full.yaml
│   │   ├── missing_project.yaml
│   │   └── old_json/
│   ├── test_loader.py
│   ├── test_schema.py
│   ├── test_adapter.py
│   ├── test_adapter_namelists.py
│   └── test_migration.py
├── test_registry/
│   ├── test_manager.py
│   ├── test_resolver.py
│   └── test_scanner_registration.py
├── test_runner/
│   └── test_local.py
└── manual/                # 手册生成器测试
    ├── conftest.py
    ├── test_escape.py
    ├── test_tables.py
    ├── test_io.py
    ├── test_generate_reference_table.py
    ├── test_generate_model_table.py
    ├── test_generate_config_schema.py
    ├── test_generate_registry_schema.py
    └── test_generate_internal_interfaces.py
\end{minted}

\section{测试金字塔}

\openbench{} 大致按经典金字塔组织：

\begin{itemize}
  \item \textbf{Unit}（基础）：单函数 / 单类（\file{test\_config/test\_schema.py}、\file{manual/test\_escape.py}）
  \item \textbf{Integration}（中层）：跨模块（\file{test\_config/test\_loader.py} 加载真 YAML、\file{test\_cli\_integration.py} 用 CliRunner 调命令）
  \item \textbf{Smoke}（顶层）：最简单的"能 import"（\file{test\_smoke.py}）
\end{itemize}

我们没有 end-to-end 测试（要跑完整 evaluation 太慢，也需大数据集）。

\section{Fixtures}

\file{tests/test\_config/fixtures/} 存放标准 YAML / JSON 输入：

\begin{itemize}
  \item \file{minimal.yaml}：最简合法配置（4 必填块 + 1 变量）
  \item \file{full.yaml}：含全部可选字段（用于回归）
  \item \file{missing\_project.yaml}：故意缺 \yamlkey{project}，用于测错误路径
  \item \file{old\_json/}：v2.0 多文件 JSON，用于测 \modname{migration}
\end{itemize}

\subsection{conftest.py 中的 pytest fixtures}

\file{tests/manual/conftest.py} 示例：

\begin{minted}{python}
@pytest.fixture
def minimal_reference_catalog(tmp_path):
    p = tmp_path / "reference_catalog.yaml"
    p.write_text(textwrap.dedent("""
        GLEAM_v4.2a_LowRes:
          name: GLEAM_v4.2a_LowRes
          ...
    """).strip())
    return p
\end{minted}

任何测试加 \texttt{minimal\_reference\_catalog} 参数自动获得该 fixture。

\section{各模块测试覆盖}

\subsection{config}

\begin{itemize}
  \item \modname{test\_schema}：dataclass 默认值、必填字段、构造器约束
  \item \modname{test\_loader}：YAML 加载、!include、错误路径、deprecation 警告
  \item \modname{test\_adapter}：dataclass → legacy 字典翻译
  \item \modname{test\_adapter\_namelists}：build\_runner\_bindings 输出的 namelist 字段完整
  \item \modname{test\_migration}：JSON / NML → YAML 转换
\end{itemize}

\subsection{registry}

\begin{itemize}
  \item \modname{test\_manager}：list/get APIs、用户级 catalog 合并
  \item \modname{test\_resolver}：reference 名解析、LowRes/MidRes 后缀
  \item \modname{test\_scanner\_registration}：scanner 注册流程
\end{itemize}

\subsection{runner}

\begin{itemize}
  \item \modname{test\_local}：run\_evaluation 主循环（用 mock NC 数据）
\end{itemize}

\subsection{cli}

\begin{itemize}
  \item \modname{test\_cli\_integration}：用 \modname{click.testing.CliRunner} 调每个命令
  \item \modname{test\_cli\_stubs}：占位命令（开发期）
\end{itemize}

\subsection{manual generators}

\begin{itemize}
  \item 5 个生成器各自有 4-5 个 test，覆盖正常输入、空输入、特殊字符转义
\end{itemize}

\section{跑测试}

\begin{minted}{bash}
# 全部
pytest

# 详细 + 单文件
pytest tests/test_config/test_loader.py -v

# 模式过滤
pytest -k "init"   # 名字含 init 的测试

# 覆盖率
pytest --cov=openbench --cov-report=term-missing

# 只跑失败的（上次失败的）
pytest --lf

# 失败立停
pytest -x

# 并发（要 pytest-xdist）
pip install pytest-xdist
pytest -n auto
\end{minted}

\section{写新测试}

\subsection{文件命名}

\begin{itemize}
  \item 路径：\file{tests/<area>/test\_<module>.py}
  \item 函数：\texttt{test\_<behavior>}（snake\_case）
  \item 类：\texttt{class Test<Subject>}（驼峰）
\end{itemize}

\subsection{TDD 风格}

新功能建议先写测试再写实现：

\begin{enumerate}
  \item 写描述行为的 test 函数（命名清楚）
  \item 跑测试，确认 fail（“红”）
  \item 写最小实现让测试通过（“绿”）
  \item Refactor，每次改完仍绿
\end{enumerate}

\subsection{测试模板}

\begin{minted}{python}
"""Test for <module>."""
import pytest

from openbench.<module> import target


def test_basic_case():
    result = target(input)
    assert result == expected


def test_edge_case_empty():
    with pytest.raises(SomeError, match="meaningful"):
        target(empty_input)


class TestComplexBehavior:
    @pytest.fixture
    def setup(self):
        return prepare_state()

    def test_path_a(self, setup):
        # ...
\end{minted}

\section{回归测试约定}

每次修复 bug 必须加回归测试。例：

\file{tests/test\_cli\_integration.py::test\_init\_output\_is\_loadable} 是因为 \cli{openbench init} 一度写出 loader 拒绝的格式而加的。即使本 bug 之后被修，测试保留以防回归。

\section{TEST\_PLAN.md / TEST\_PLAN\_CLI.md / TEST\_PLAN\_GUI.md}

仓库根有三份 TEST\_PLAN 文档，列出"应该测什么"。这些是规划，不是自动跑的。每开发 milestone，更新这三份文档对照实际覆盖。

\warninline{TEST\_PLAN 可能过期（每个 milestone 末尾刷新）。如果你看到与代码不符，请提 PR 更新或在 GitHub Issue 标记。}

\section{CI 期望}

\openbench{} 的 GitHub Actions 应至少跑：

\begin{itemize}
  \item Python 3.10 / 3.11 / 3.12 矩阵
  \item ruff check + format check
  \item pytest（带覆盖率）
\end{itemize}

GUI / SSH / 重 IO 测试不在 CI（环境复杂）。

\section{测试覆盖率目标}

\begin{itemize}
  \item 新代码：> 80\%
  \item config / runner / registry：~90\%（基础设施类）
  \item visualization：~50\%（图形难自动测）
  \item gui：当前未测
\end{itemize}

\section{下一步}

最后一章：贡献流程（commit、PR、release）。
